% Auto-generated complete chapter from 28D_artstation_portfolio_benchmark_complete.md
\chapter{28D — Benchmark actualizado de portafolios técnicos y arte 3D en tiempo real}
\label{complete:root_036}

\begin{sourcemeta}
\textbf{Fuente:} \href{file:///E:/Laboral/28D_artstation_portfolio_benchmark_complete.md}{28D\_artstation\_portfolio\_benchmark\_complete.md}\\
\textbf{Seccion:} root\\
\textbf{Estado:} canonical\\
\textbf{Rol:} Version consolidada\\
\textbf{Lineas originales:} 855\\
\textbf{Ultima modificacion:} 2026-06-11 18:19\\
\textbf{Deduplicacion conservadora:} 0 bloques repetidos omitidos; 0 caracteres repetidos no reimpresos.
\end{sourcemeta}

\section*{Resumen de orientacion}
--- module: 28D\_artstation\_portfolio\_benchmark status: updated\_complete source\_task: A4\_retry\_accessible\_portfolio\_benchmark scope: accessible\_portfolios\_and\_technical\_breakdowns artstation\_direct\_access: blocked\_by\_cloudflare date: 2026-06-11 owner\_profile: Alexander Woodcock Salomon target\_positioning: - Real-Time 3D Developer - Unity Technical Artist - Unity WebGL Developer - Technical Visualization Developer - XR Developer - Simulation Developer - Tools / Pipeline Develop

\section*{Contenido completo depurado}
\addcontentsline{toc}{section}{Contenido completo depurado}

\medskip\hrule\medskip


module: 28D\_artstation\_portfolio\_benchmark

status: updated\_complete

source\_task: A4\_retry\_accessible\_portfolio\_benchmark

scope: accessible\_portfolios\_and\_technical\_breakdowns

artstation\_direct\_access: blocked\_by\_cloudflare

date: 2026-06-11

owner\_profile: Alexander Woodcock Salomon

target\_positioning:


\begin{itemize}
\item Real-Time 3D Developer
\item Unity Technical Artist
\item Unity WebGL Developer
\item Technical Visualization Developer
\item XR Developer
\item Simulation Developer
\item Tools / Pipeline Developer
\end{itemize}


\medskip\hrule\medskip







\subsection{1. Estado del módulo}




Este módulo reemplaza la versión anterior de \texttt{28D\_artstation\_portfolio\_benchmark.md}.




La tarea original buscaba un benchmark centrado en ArtStation. Sin embargo, el acceso directo a ArtStation quedó bloqueado por verificación de Cloudflare. Por tanto, este módulo no debe presentarse como un benchmark exhaustivo de ArtStation.




El módulo actualizado usa fuentes públicas accesibles equivalentes: portafolios personales, GitHub Pages, Framer, blogs técnicos y breakdowns públicos de arte técnico, shaders, VFX, herramientas y 3D en tiempo real.




\subsection{2. Decisión metodológica}




ArtStation sigue siendo una fuente valiosa para evaluar presentación visual. No obstante, para estudiar estructura técnica, storytelling de proyecto, documentación, código, métricas y breakdowns, las fuentes accesibles recopiladas son más útiles que perfiles visuales bloqueados.




La decisión correcta es separar:




\begin{itemize}
\item \textbf{ArtStation directo:} pendiente por bloqueo de acceso.
\item \textbf{Benchmark visual/técnico accesible:} válido para patrones de portafolio.
\item \textbf{Case studies técnicos:} fuente principal para construir los proyectos de Alexander.
\item \textbf{GitHub/README:} fuente principal para validar implementación y código.
\end{itemize}




\subsection{3. Calidad del benchmark A4\_retry}




\begin{center}
\scriptsize
\begin{longtable}{>{\raggedright\arraybackslash}p{0.455\linewidth} >{\raggedright\arraybackslash}p{0.455\linewidth}}
\toprule
Criterio & Resultado \\
\midrule
\endfirsthead
\toprule
Criterio & Resultado \\
\midrule
\endhead
Fuentes accesibles & Sí \\
ArtStation directo & No \\
Portafolios técnicos & Sí \\
Blogs técnicos & Sí \\
Shaders/VFX & Sí \\
Unity/Unreal & Sí \\
Métricas & Escasas \\
Código visible & Parcial \\
Utilidad para Alexander & Alta \\
Exhaustividad & Media \\
\bottomrule
\end{longtable}
\end{center}




\subsection{4. Inventario de fuentes válidas}




\begin{center}
\scriptsize
\begin{longtable}{>{\raggedright\arraybackslash}p{0.225\linewidth} >{\raggedright\arraybackslash}p{0.225\linewidth} >{\raggedright\arraybackslash}p{0.225\linewidth} >{\raggedright\arraybackslash}p{0.225\linewidth}}
\toprule
Fuente & Plataforma & Foco & Calidad \\
\midrule
\endfirsthead
\toprule
Fuente & Plataforma & Foco & Calidad \\
\midrule
\endhead
Alla Eddine Rakik & Framer & Tech art general & Media \\
Rakik — Boom Mike VR & Framer & VR demo & Baja-media \\
Rakik — Shader VFX Pack & Framer & Shader/VFX & Baja-media \\
Jettelly — Potion Shader & Blog & Unity shader breakdown & Alta \\
Harry Alisavakis & Blog & Technical art resources & Media \\
Omar Rinaz Costa & GitHub Pages & CV + tech art & Media \\
Alson Entuna & Personal site & Tools/pipeline & Media \\
Wolf van Veen & Personal site & Shaders/VFX & Alta \\
Francisco Múrias & GitHub Pages & VFX/rendering & Alta \\
Julhe — SDF Textures & Blog & Shader technique & Alta \\
Julhe — Shader Optimization & Blog & Shader performance & Alta \\
Julhe — X5 Shader & Blog & Shader framework & Alta \\
\bottomrule
\end{longtable}
\end{center}




\subsection{5. Fuentes fuertes}




\subsubsection{5.1 Jettelly — Real-Time Potion Shader Breakdown}




URL: \url{https://jettelly.com/blog/a-real-time-potion-shader-breakdown-in-unity}




Es uno de los ejemplos más útiles porque no se limita a mostrar una imagen final. Presenta el problema técnico, la solución visual, el uso de Unity URP, Shader Graph y C\#, y el comportamiento del efecto en movimiento.




Patrones adaptables:




\begin{itemize}
\item Título específico.
\item Resultado visual temprano.
\item Explicación de comportamiento.
\item Capturas intermedias.
\item Video o GIF del resultado.
\item Separación entre efecto visual y técnica usada.
\end{itemize}




Aplicación para Alexander:




Este formato sirve para un breakdown de:




\begin{itemize}
\item modo X-Ray;
\item shader de clipping/cross-section;
\item visualización térmica heurística;
\item resaltado de piezas;
\item sistema de selección;
\item controles WebGL.
\end{itemize}




Lo que no debe copiar:




No basta con decir “hecho en Unity”. El valor está en explicar la relación entre intención visual, limitaciones del motor, implementación y resultado.




\subsubsection{5.2 Julhe — Always-Sharp SDF Textures}




URL: \url{https://julhe.github.io/blog/always-sharp-sdf-textures/}




Este ejemplo funciona como referencia de comunicación técnica avanzada. Explica una técnica concreta, desarrolla el problema y muestra razonamiento técnico.




Patrones adaptables:




\begin{itemize}
\item Problema técnico acotado.
\item Explicación conceptual.
\item Visuales comparativos.
\item Código o pseudocódigo.
\item Resultado final verificable.
\end{itemize}




Aplicación para Alexander:




Sirve como modelo para escribir piezas técnicas sobre:




\begin{itemize}
\item optimización de lectura visual en WebGL;
\item diseño de hotspots;
\item claridad de bordes;
\item visualización de piezas pequeñas;
\item control de transparencia;
\item UI técnica sobre modelos complejos.
\end{itemize}




Lo que no debe copiar:




No conviene publicar artículos demasiado matemáticos si el objetivo principal es empleabilidad inicial. La profundidad técnica debe ser suficiente, pero legible para recruiters técnicos y leads.




\subsubsection{5.3 Julhe — How (Not) To Optimize Shaders}




URL: \url{https://julhe.github.io/blog/how-not-to-optimize-shaders/}




Este ejemplo es especialmente relevante porque enseña mediante anti-patrones. No solo muestra qué hacer, sino qué evitar.




Patrones adaptables:




\begin{itemize}
\item Planteamiento de error común.
\item Comparación entre enfoques.
\item Análisis de rendimiento.
\item Recomendaciones prácticas.
\item Conclusión aplicable.
\end{itemize}




Aplicación para Alexander:




Puede inspirar un artículo propio:




\begin{itemize}
\item “How not to optimize a Unity WebGL technical viewer”.
\item “Lessons from building a browser-based drone assembly viewer”.
\item “What I would rebuild in my thesis viewer”.
\item “Mobile-first constraints in Unity WebGL visualization”.
\end{itemize}




Lo que no debe copiar:




No debe convertirse en una crítica abstracta. Debe estar conectado a una experiencia propia y a decisiones reales del proyecto.




\subsubsection{5.4 Francisco Múrias — Technical Art Portfolio}




URL: \url{https://franciscomurias.github.io/}




Este portafolio es útil por su grilla visual, foco en proyectos técnicos, y combinación de arte técnico con rendering, VFX y herramientas.




Patrones adaptables:




\begin{itemize}
\item Home visual fuerte.
\item Proyectos en tarjetas.
\item Separación entre trabajos personales y profesionales.
\item Enlaces a GitHub cuando aplica.
\item Estética técnica sin sobreexplicar en la portada.
\end{itemize}




Aplicación para Alexander:




Sirve para estructurar una landing de portafolio con:




\begin{itemize}
\item TwinSight X500 como proyecto principal;
\item miniaturas de módulos técnicos;
\item tarjetas de herramientas;
\item sección de breakdowns;
\item enlaces a demo y GitHub.
\end{itemize}




Lo que no debe copiar:




No debe depender solo de imágenes atractivas. En el caso de Alexander, el valor competitivo está en el proceso técnico y la interacción WebGL.




\subsubsection{5.5 Wolf van Veen — Technical Artist}




URL: \url{https://wolfvanveen.art/}




Este ejemplo destaca por la organización visual mediante categorías: post-processing, shaders, particles, models y textures. Aunque el propio sitio indica que el portafolio está desactualizado, su estructura sigue siendo útil.




Patrones adaptables:




\begin{itemize}
\item Categorías visibles.
\item Navegación simple.
\item Galerías por proyecto.
\item Modal o página individual por pieza.
\item Herramientas mostradas por proyecto.
\end{itemize}




Aplicación para Alexander:




El portafolio de Alexander puede usar categorías similares, adaptadas a su perfil:




\begin{itemize}
\item Technical Visualization.
\item Unity WebGL.
\item Shaders \& View Modes.
\item Tools \& Pipeline.
\item 3D Assets.
\item XR / Interactive.
\item Case Studies.
\end{itemize}




Lo que no debe copiar:




No debe dejar proyectos sin explicación técnica. Cada tarjeta debe responder qué problema resolvió y cómo se implementó.




\subsection{6. Fuentes medias}




\subsubsection{6.1 Alla Eddine Rakik — Technical Art Portfolio}




URL: \url{https://technical-art-portfolio.io.framer.ai/}




El sitio tiene una estructura clara: presentación, herramientas, lenguajes y grilla de proyectos. Es útil como referencia visual de organización.




Aporta:




\begin{itemize}
\item landing directa;
\item lista de software;
\item proyectos por miniaturas;
\item variedad de áreas: VR, 3D, VFX, OpenGL.
\end{itemize}




Limitación:




Los proyectos tienen poca profundidad técnica. Sirven como muestra visual, pero no como case studies completos.




Aplicación para Alexander:




Puede servir como modelo inicial de landing, pero debe complementarse con páginas de proyecto mucho más documentadas.




\subsubsection{6.2 Omar Rinaz Costa — Portfolio}




URL: \url{https://omarrinaz.github.io/}




Este portafolio sirve para observar estructura curricular: about, experiencia, educación, skills y portfolio. Es útil para la transición entre CV y portafolio.




Aporta:




\begin{itemize}
\item línea de tiempo;
\item roles y fechas;
\item educación visible;
\item lista de responsabilidades;
\item botón de descarga de CV.
\end{itemize}




Limitación:




Faltan métricas y proyectos técnicos profundos.




Aplicación para Alexander:




Puede inspirar una página \texttt{About / CV}, pero no debe ser el modelo principal para proyectos técnicos.




\subsubsection{6.3 Alson Entuna — Lead Technical Artist Pipeline \& Tools}




URL: \url{https://alsonentuna.github.io/}




El sitio es útil porque comunica una transición profesional: gameplay programming hacia technical art, tools y pipeline.




Aporta:




\begin{itemize}
\item posicionamiento claro;
\item foco en tools/pipeline;
\item relato profesional;
\item seniority explícito.
\end{itemize}




Limitación:




El home muestra más perfil que evidencia técnica. Para Alexander, esto sería insuficiente sin proyectos concretos.




Aplicación para Alexander:




Puede inspirar el texto de posicionamiento profesional:




\begin{quote}
Multimedia engineer focused on real-time 3D visualization, Unity WebGL, technical art systems and interactive product visualization.
\end{quote}




\subsubsection{6.4 Harry Alisavakis — Technical Art Blog}




URL: \url{https://halisavakis.com/}




El blog funciona como repositorio de recursos y artículos. Es útil como referencia de consistencia y marca personal.




Aporta:




\begin{itemize}
\item publicaciones recurrentes;
\item recursos técnicos;
\item enfoque educativo;
\item mezcla Unity/Unreal;
\item comunidad técnica.
\end{itemize}




Limitación:




Es más blog que portafolio. No todo está organizado para evaluación rápida de recruiter.




Aplicación para Alexander:




Sirve como modelo para una sección \texttt{Technical Notes} o \texttt{Breakdowns}, no como home principal.




\subsection{7. Fuentes débiles o incompletas}




\subsubsection{7.1 Rakik — Boom Mike VR}




URL: \url{https://technical-art-portfolio.io.framer.ai/boom-mike}




Aporta una demo VR con video, pero poca explicación técnica. Es útil como referencia mínima de presentación visual.




No debe ser modelo principal para Alexander.




\subsubsection{7.2 Rakik — Shader VFX Pack}




URL: \url{https://technical-art-portfolio.io.framer.ai/shader-vfx-pack}




Aporta una muestra visual de efectos, pero no explica implementación, problema, parámetros ni resultados.




Debe evitarse este patrón si el objetivo es aplicar a roles técnicos.




\subsection{8. Patrones de estructura para el portafolio de Alexander}




\subsubsection{8.1 Landing page}




La landing debe responder en menos de 10 segundos:




\begin{itemize}
\item quién es Alexander;
\item qué rol busca;
\item cuál es su proyecto principal;
\item qué evidencia técnica puede revisar el evaluador;
\item dónde están demo, GitHub, CV y contacto.
\end{itemize}




Estructura recomendada:




\begin{mdcode}
Bloque de codigo: text
Hero
└── Role headline
└── One-line positioning
└── CTA: View TwinSight X500
└── CTA: GitHub
└── CTA: Download CV

Featured Case Study
└── TwinSight X500
└── Unity WebGL technical viewer
└── Demo / GitHub / Breakdown

Project Categories
└── Technical Visualization
└── Unity WebGL
└── Shaders & View Modes
└── Tools & Pipeline
└── 3D Assets
└── Technical Notes
\end{mdcode}




\subsubsection{8.2 Project card}




Cada tarjeta de proyecto debe mostrar:




\begin{itemize}
\item nombre;
\item rol;
\item tecnología;
\item tipo de problema;
\item resultado visible;
\item enlace a demo;
\item enlace a breakdown;
\item enlace a código si aplica.
\end{itemize}




Formato sugerido:




\begin{mdcode}
Bloque de codigo: markdown
## TwinSight X500

**Role:** Unity WebGL Developer / Technical Artist  
**Type:** Browser-based 3D technical visualization  
**Tools:** Unity 6000, Blender, WebGL, URP, GitHub Pages  
**Focus:** Spatial inspection, part selection, exploded view, cross-section, heuristic thermal reading  
**Demo:** [link]  
**Code:** [link]  
**Case Study:** [link]
\end{mdcode}




\subsubsection{8.3 Case study técnico}




La estructura debe tomar como referencia Jettelly y Julhe.




Plantilla:




\begin{mdcode}
Bloque de codigo: markdown
# Project name

## Context
What problem existed?

## Goal
What the system needed to solve.

## Constraints
Browser, mobile, WebGL, asset size, usability, academic scope.

## Technical approach
Architecture, data, shaders, UI, pipeline.

## Implementation
Key systems and decisions.

## Results
Screenshots, video, metrics, validation.

## What I learned
Trade-offs, mistakes, next iteration.

## Links
Demo, GitHub, report, video.
\end{mdcode}




\subsection{9. Qué debe adaptar Alexander}




\subsubsection{9.1 De Jettelly}




Adaptar:




\begin{itemize}
\item breakdown visual paso a paso;
\item explicación de comportamiento dinámico;
\item uso de imágenes y video;
\item foco en Unity.
\end{itemize}




Aplicar a:




\begin{itemize}
\item thermal heuristic view;
\item X-Ray mode;
\item cross-section shader;
\item exploded view;
\item selection highlighting.
\end{itemize}




\subsubsection{9.2 De Julhe}




Adaptar:




\begin{itemize}
\item profundidad técnica;
\item explicación de problema;
\item anti-patrones;
\item código o pseudocódigo;
\item análisis de rendimiento.
\end{itemize}




Aplicar a:




\begin{itemize}
\item WebGL optimization;
\item shader modes;
\item UI performance;
\item model hierarchy normalization;
\item asset import pipeline.
\end{itemize}




\subsubsection{9.3 De Francisco Múrias}




Adaptar:




\begin{itemize}
\item grilla de proyectos;
\item estética técnica;
\item tarjetas visuales;
\item enlaces a GitHub;
\item separación por tipo de trabajo.
\end{itemize}




Aplicar a:




\begin{itemize}
\item home del portafolio;
\item featured projects;
\item technical demos;
\item visual archive.
\end{itemize}




\subsubsection{9.4 De Wolf van Veen}




Adaptar:




\begin{itemize}
\item categorías claras;
\item filtros por tipo de trabajo;
\item modal o galería de proyecto;
\item herramientas visibles por pieza.
\end{itemize}




Aplicar a:




\begin{itemize}
\item portfolio navigation;
\item project taxonomy;
\item visual browsing.
\end{itemize}




\subsubsection{9.5 De Omar Rinaz}




Adaptar:




\begin{itemize}
\item timeline de experiencia;
\item formación académica;
\item lista de roles;
\item sección CV.
\end{itemize}




Aplicar a:




\begin{itemize}
\item About;
\item CV page;
\item career context.
\end{itemize}




\subsection{10. Qué no debe copiar Alexander}




Evitar:




\begin{itemize}
\item portafolios con solo imágenes;
\item proyectos sin contexto;
\item ausencia de métricas;
\item falta de enlaces a demo;
\item listas largas de software sin evidencia;
\item describirse como senior sin respaldo;
\item claims no verificables;
\item usar ArtStation como única superficie;
\item esconder el GitHub;
\item omitir decisiones técnicas.
\end{itemize}




\subsection{11. Recomendaciones concretas para el portafolio}




\subsubsection{11.1 Proyecto principal}




TwinSight X500 debe ser el proyecto hero.




Título recomendado:




\begin{mdcode}
Bloque de codigo: text
TwinSight X500 — Unity WebGL Technical Visualization Viewer
\end{mdcode}




Subtítulo recomendado:




\begin{mdcode}
Bloque de codigo: text
Browser-based interactive 3D viewer for inspecting the Holybro X500 V2 drone assembly, focused on spatial understanding, part hierarchy, exploded view, cross-section and heuristic visual analysis.
\end{mdcode}




\subsubsection{11.2 Categorías}




Usar estas categorías:




\begin{mdcode}
Bloque de codigo: text
Technical Visualization
Unity WebGL
Shaders & View Modes
Tools & Pipeline
3D Production
Academic Research
Technical Notes
\end{mdcode}




\subsubsection{11.3 Case studies mínimos}




Publicar al menos 4 piezas:




\begin{mdcode}
Bloque de codigo: text
1. TwinSight X500 — full case study
2. Cross-section shader / clipping plane
3. Drone CAD-to-Web pipeline
4. WebGL optimization / mobile constraints
\end{mdcode}




Deseables:




\begin{mdcode}
Bloque de codigo: text
5. Fastener modular system
6. Exploded view logic
7. UI bottom-sheet interaction system
8. Heuristic thermal visualization
\end{mdcode}




\subsubsection{11.4 Medios visuales}




Incluir:




\begin{itemize}
\item hero video corto;
\item GIF de selección de piezas;
\item GIF de exploded view;
\item GIF de cross-section;
\item capturas before/after;
\item diagrama de pipeline;
\item diagrama de arquitectura;
\item tabla de métricas;
\item enlace a demo.
\end{itemize}




\subsubsection{11.5 Métricas}




Agregar cuando sean verificables:




\begin{mdcode}
Bloque de codigo: text
FPS
frame time
build size
load time
memory
draw calls
triangles
texture memory
SUS
NASA-TLX Raw
task completion
\end{mdcode}




No inventar métricas. Si no están medidas, marcar como pendiente.




\subsection{12. Cambios concretos para módulos previos}




\subsubsection{12.1 Actualizar \texttt{17\_cv\_base\_and\_role\_variants.md}}




Agregar una sección de proyectos técnicos con enfoque medible:




\begin{mdcode}
Bloque de codigo: markdown
### Featured Technical Project

**TwinSight X500 — Unity WebGL Technical Visualization Viewer**  
Built an interactive browser-based 3D viewer for inspecting the Holybro X500 V2 drone assembly. Implemented part hierarchy, selection, exploded view, cross-section and visual analysis modes for technical study and usability evaluation.
\end{mdcode}




Cuando existan métricas finales, añadir:




\begin{mdcode}
Bloque de codigo: markdown
Measured WebGL performance across load time, frame time and memory usage.
\end{mdcode}




\subsubsection{12.2 Actualizar README del repositorio}




El README ya está bien orientado, pero puede mejorar como portafolio técnico si añade:




\begin{itemize}
\item GIF hero;
\item project architecture diagram;
\item pipeline diagram;
\item screenshots table;
\item direct link to case study;
\item measured performance section;
\item “What I built” section.
\end{itemize}




\subsubsection{12.3 Actualizar landing de GitHub Pages}




La landing debe priorizar:




\begin{enumerate}
\item Demo inmediata.
\item Caso de estudio.
\item Capturas del flujo.
\item Métricas verificadas.
\item Enlaces a manuales.
\item CV/portfolio.
\end{enumerate}




\subsection{13. Estructura recomendada del sitio personal}




\begin{mdcode}
Bloque de codigo: text
/
├── Home
├── Projects
│   ├── TwinSight X500
│   ├── Cross-section Shader
│   ├── CAD-to-Web Pipeline
│   └── WebGL Optimization Notes
├── Technical Notes
│   ├── Unity WebGL constraints
│   ├── Shader view modes
│   └── 3D model optimization
├── About
├── CV
└── Contact
\end{mdcode}




\subsection{14. Taxonomía de proyectos}




\begin{center}
\scriptsize
\begin{longtable}{>{\raggedright\arraybackslash}p{0.455\linewidth} >{\raggedright\arraybackslash}p{0.455\linewidth}}
\toprule
Categoría & Uso \\
\midrule
\endfirsthead
\toprule
Categoría & Uso \\
\midrule
\endhead
Technical Visualization & TwinSight, drone viewer \\
Unity WebGL & Browser deployment \\
Shaders & X-Ray, thermal, clipping \\
Tools & pipeline utilities \\
3D Production & Blender, bake, optimization \\
Research & thesis validation \\
Notes & technical articles \\
\bottomrule
\end{longtable}
\end{center}




\subsection{15. Plantilla para página de proyecto}




\begin{mdcode}
Bloque de codigo: markdown
---
title: TwinSight X500
role: Unity WebGL Developer / Technical Artist
tools: Unity, Blender, WebGL, URP
status: public demo
year: 2026
---

# TwinSight X500

## Summary

TwinSight X500 is a browser-based interactive 3D viewer for technical inspection of the Holybro X500 V2 drone assembly.

## Problem

Traditional 2D manuals force users to mentally reconstruct spatial relationships between parts.

## Solution

The project provides an interactive 3D environment with selection, part hierarchy, exploded view, cross-section and visual analysis modes.

## My Role

I designed and implemented the Unity WebGL prototype, prepared the 3D asset pipeline and developed the interaction systems used in the final thesis demo.

## Technical Stack

- Unity 6000
- WebGL
- URP
- Blender
- GitHub Pages
- C#

## Key Systems

- App state machine
- Part selection
- Bottom-sheet UI
- Exploded view
- Cross-section
- View modes
- Thermal heuristic visualization

## Constraints

- Browser deployment
- Mobile-first interaction
- Academic validation scope
- Complex imported hierarchy
- Performance constraints

## Results

Add measured metrics here when available.

## Media

Add GIFs, screenshots and video.

## Links

- Demo
- GitHub
- Report
- Technical manual
\end{mdcode}




\subsection{16. Plantilla para breakdown técnico}




\begin{mdcode}
Bloque de codigo: markdown
# Cross-section shader breakdown

## What I wanted to solve

Users needed to inspect internal spatial relationships without losing context of the full drone assembly.

## Constraint

The solution had to work in Unity WebGL and remain usable on mobile devices.

## Approach

I implemented a shader-based clipping plane controlled through the Analyze mode UI.

## Implementation

Explain the clipping logic, material setup, UI control and limitations.

## Result

Show GIF, screenshots and performance notes.

## Lessons

Explain what worked, what failed and what would be improved.
\end{mdcode}




\subsection{17. Riesgos}




\begin{center}
\scriptsize
\begin{longtable}{>{\raggedright\arraybackslash}p{0.455\linewidth} >{\raggedright\arraybackslash}p{0.455\linewidth}}
\toprule
Riesgo & Mitigación \\
\midrule
\endfirsthead
\toprule
Riesgo & Mitigación \\
\midrule
\endhead
Portafolio solo visual & Añadir breakdowns \\
Claims sin evidencia & Medir antes de publicar \\
GitHub oculto & Enlace visible \\
Proyecto académico percibido como menor & Enfatizar alcance técnico \\
Falta de seniority & Posicionarse como junior/mid técnico \\
Demasiado texto & Usar jerarquía visual \\
Falta de demo & Priorizar build estable \\
\bottomrule
\end{longtable}
\end{center}




\subsection{18. Orden operativo recomendado}




\begin{mdcode}
Bloque de codigo: text
1. Congelar demo pública.
2. Capturar GIFs clave.
3. Medir rendimiento verificable.
4. Crear case study principal.
5. Crear 2–3 breakdowns técnicos.
6. Actualizar README.
7. Construir landing personal.
8. Ajustar CV.
9. Publicar LinkedIn/GitHub profile.
\end{mdcode}




\subsection{19. Prompt para A5}




\begin{mdcode}
Bloque de codigo: text
Execute A5_case_study_benchmark.

Objective:
Analyze 3–5 strong accessible technical case studies relevant to Alexander’s positioning:
Unity Technical Artist, Real-Time 3D Developer, Unity WebGL Developer, Technical Visualization Developer, Shader/Tools Developer.

Prioritize:
- Jettelly Real-Time Potion Shader Breakdown;
- Julhe shader breakdowns;
- Francisco Múrias technical projects;
- other accessible technical art breakdowns with code, images, video or metrics.

Output:
- source inventory;
- case study structure;
- storytelling pattern;
- technical depth;
- visual evidence;
- metrics used;
- what Alexander should adapt;
- project page template;
- breakdown article template.

Rules:
- Use public accessible sources only.
- Do not use ArtStation if blocked.
- Do not invent missing data.
- Keep tables concise.
- Provide actionable recommendations.
\end{mdcode}




\subsection{20. Conclusión}




El benchmark actualizado muestra que el portafolio de Alexander no debe depender de una galería tipo ArtStation. Su ventaja competitiva está en combinar:




\begin{itemize}
\item evidencia visual;
\item explicación técnica;
\item demo WebGL;
\item código público;
\item métricas verificables;
\item claridad de rol;
\item estudios de caso.
\end{itemize}




La prioridad debe ser construir un portafolio híbrido: visual como un artista técnico, documentado como un desarrollador, y medible como un proyecto de ingeniería.



